\section{Tools and Environment}\label{sec:tools}

Table~\ref{tab:tools} lists the tools and libraries with the versions used. The versions of the development tools were
read from the tools themselves (\texttt{java -version}, \texttt{git --version}, \texttt{code --version}); those of the
libraries come from the project's \texttt{pom.xml} and from the versions Maven resolved for it.

\begin{table}[!t]
  \caption{Tools and libraries}\label{tab:tools}
  \centering\footnotesize
  \begin{tabularx}{\linewidth}{@{}L{0.33\linewidth}L{0.17\linewidth}Y@{}}
    \toprule
    Tool or library & Version & Used for \\
    \midrule
    Windows 11 Pro & 10.0.26300 & Operating system \\
    Oracle JDK & 25.0.1 LTS & Compiling, running and debugging; the \texttt{javadoc} tool \\
    Visual Studio Code & 1.147.0 & IDE \\
    Extension Pack for Java & 0.31.2 & Java support in the IDE (below) \\
    Apache Maven, run by the Maven wrapper 3.3.4 & 3.10.0 & Build, dependencies, tests, Javadoc \\
    Git for Windows & 2.54.0 & Version control \\
    GitHub & --- & Remote repository \\
    Spring Boot & 4.1.3 & Dependency injection, transactions, configuration \\
    JavaFX & 27.0.1 & Graphical user interface (FXML and CSS) \\
    Hibernate ORM & 7.4.15 & Mapping the task onto its table \\
    SQLite JDBC driver & 3.53.4.0 & The database file \\
    Flyway & 12.4.0 & Database schema migrations \\
    JUnit Jupiter & 6.0.3 & Automated tests \\
    Maven Javadoc Plugin & 3.12.0 & Generating the HTML documentation \\
    LaTeX (IEEEtran), Tectonic & 0.15.0 & Typesetting this report \\
    \bottomrule
  \end{tabularx}
\end{table}

\subsection{The IDE}

The brief names NetBeans, IntelliJ IDEA and Eclipse as example IDEs. We used Visual Studio Code with Microsoft's
Extension Pack for Java~\cite{vscode-java}. Its language support (Language Support for Java by Red Hat, 1.58.0) is
built on the Eclipse JDT Language Server, so the code is compiled and checked as it is typed by the same Java tooling
Eclipse uses. The pack also brings the Debugger for Java (0.60.0) for breakpoints, stepping and inspecting variables,
the Test Runner for Java (0.47.0), Maven for Java (0.46.0), and the Project Manager for Java (0.27.7), whose Java
Projects view shows the project's source folders, dependencies and JDK. How the project was set up and run
in the IDE is described in Section~\ref{sec:ide}.

\subsection{The JDK and the build}

The project targets Java~25 (\texttt{java.version} in \texttt{pom.xml}) and was built and run on Oracle JDK~25.0.1,
with \texttt{JAVA\_HOME} pointing at it. It is a Maven project with the Maven wrapper checked in: \texttt{./mvnw}
downloads Maven~3.10.0 on first use, so the project builds the same way from the IDE, from a terminal, and on another
computer that has only the JDK. The commands used throughout are:
\begin{itemize}
  \item \texttt{./mvnw spring-boot:run}: start the application;
  \item \texttt{./mvnw test}: run the JUnit tests;
  \item \texttt{./mvnw javadoc:javadoc}: generate the HTML documentation.
\end{itemize}
JavaFX and the SQLite driver both load native libraries, which Java~25 warns about unless native access is allowed.
The \texttt{pom.xml} therefore passes \texttt{--enable-native-access=ALL-UNNAMED} to the application when Maven starts
it and to the tests, and writes the same permission into the jar's manifest; the IDE's launch configuration passes
it too.
